\reading{CR-13}{Credit Derivatives}{DEFER}{2}

\begin{crkeybox}[Learning objectives for this reading]
\begin{enumerate}[label=\textbf{\alph*.},leftmargin=2em]
\item Describe a credit derivative, credit default swap (CDS), total return swap, and collateralized debt obligation (CDO).
\item Explain how to account for credit risk exposure in valuing a CDS.
\item Identify the default probabilities used to value a CDS.
\item Evaluate the use of credit indices and fixed coupons in pricing CDS transactions.
\item Define CDS forwards and CDS options.
\item Describe the process of valuing a synthetic CDO using the spread payments approach and the Gaussian copula model of time to default approach.
\item Define the two measures of implied correlation: compound (tranche) correlation and base correlation.
\item Discuss alternative approaches used to estimate default correlation.
\end{enumerate}
\end{crkeybox}

\begin{crtrapbox}[Single-layer reading]
CR-13 exists only in the Crash layer.
\end{crtrapbox}

% =====================================================================
\lo{CR-13 a}{Describe a credit derivative, credit default swap (CDS), total return swap, and collateralized debt obligation (CDO).}

\subsection*{Credit default swap}

A CDS is insurance against the risk of default by a specific company, the
\term{reference entity}. It gives the buyer the right to sell bonds at face
value, and imposes on the seller the obligation to buy bonds at face value, when
a credit event occurs. The full mechanics are at \textbf{CR-9 g}.

\begin{itemize}
\item \term{CDS spread.} The total amount paid per year, as a percentage of
notional principal, to buy protection.
\item \term{Credit event definition.} Failure to make a payment, debt
restructuring, or bankruptcy.
\item \term{Settlement methods.} \emph{Physical settlement}: the buyer sells
bonds to the seller at face value. \emph{Cash settlement}: a two-stage auction
process determines the mid-market value of the cheapest deliverable bond, and the
cash payoff is based on the auction results.
\end{itemize}

\gapbadge

{\footnotesize\color{FRMGrey}The objective names four things. The source notes
describe the CDS and, at objective f, the synthetic CDO, but never define a total
return swap or a CDO. The following is written from the GARP curriculum.}

\begin{crgapbox}[What the objective asks for: total return swap and CDO]
\textbf{Total return swap.} A total return swap is a type of credit derivative.
It is an agreement to exchange the \emph{total return} on a bond, or any
portfolio of assets, for a floating rate plus a spread. The total return includes
coupons, interest, and the gain or loss on the asset over the life of the swap.

\smallskip
An example is a 5-year agreement with a notional principal of \$100 million to
exchange the total return on a corporate bond for the floating rate plus 25 basis
points. On coupon payment dates the \emph{payer} pays the coupons earned on an
investment of \$100 million in the bond; the \emph{receiver} pays interest at the
floating rate plus 25 basis points on a principal of \$100 million. At the end of
the life of the swap there is a payment reflecting the change in value of the
bond: if the bond rises 10\%, the payer pays \$10 million; if it falls 15\%, the
receiver pays \$15 million. If there is a default on the bond, the swap is
usually terminated and the receiver makes a final payment equal to the excess of
\$100 million over the market value of the bond.

\smallskip
\textbf{Collateralized debt obligation.} An asset-backed security where the
underlying assets are \emph{bonds} is known as a collateralized debt obligation.
A waterfall is defined for the interest and principal payments on the bonds. The
precise rules underlying the waterfall are complicated, but they are designed to
ensure that if one tranche is more senior than another, it is more likely to
receive promised interest payments and repayments of principal.

\smallskip
When a CDO is created from a bond portfolio in this way, the resulting structure
is a \term{cash CDO}, as distinct from the synthetic CDO at objective f.
\end{crgapbox}

% =====================================================================
\lo{CR-13 b}{Explain how to account for credit risk exposure in valuing a CDS.}

Valuing a CDS means equating two present values: what the protection buyer expects
to pay, and what the protection buyer expects to receive.

\begin{crfmlbox}[The two legs, per period]
\[ \text{expected payment} \;=\; \text{probability of survival} \times
   \text{notional principal} \times s \]
\[ \text{expected payoff} \;=\; (1-\text{recovery rate}) \times \mathrm{PD}
   \times \text{notional principal} \]
\vk{$s$ = the CDS spread. Payments are made only if the entity has \emph{survived}
to the payment date, so they are weighted by the survival probability; payoffs
occur only on default, so they are weighted by the default probability.}
\end{crfmlbox}

Defaults are assumed to occur \emph{mid-period}. That creates a third leg: if
default occurs halfway through a year, the protection buyer owes half a year of
accrued premium.

\begin{center}
\begin{tikzpicture}[font=\sffamily\footnotesize]
\draw[FRMGrey, line width=1pt] (0,0) -- (11.4,0);
\foreach \x/\lab in {0/0, 3.6/1, 7.2/2, 10.8/3}{
  \draw[FRMGrey, line width=0.8pt] (\x,-0.12) -- (\x,0.12);
  \node[below=3pt, font=\scriptsize] at (\x,-0.12) {\lab};
}
% payments up at 1,2,3
\foreach \x in {3.6,7.2,10.8}
  \draw[-{Stealth[length=5pt]}, FRMNavy, line width=1.1pt] (\x,0.1) -- (\x,1.35);
\node[font=\scriptsize\bfseries, text=FRMNavy, anchor=west] at (0.15,1.35)
  {expected payments (survival $\times$ notional $\times s$)};
% payoffs down at 0.5,1.5,2.5
\foreach \x in {1.8,5.4,9.0}
  \draw[-{Stealth[length=5pt]}, FRMGold, line width=1.1pt] (\x,-0.1) -- (\x,-1.35);
\node[font=\scriptsize\bfseries, text=FRMGold, anchor=west] at (0.15,-1.6)
  {expected payoffs $(1-R)\times\mathrm{PD}\times$ notional, and accrual $=\mathrm{PD}\times s/2$};
\foreach \x/\lab in {1.8/0.5, 5.4/1.5, 9.0/2.5}
  \node[font=\scriptsize, text=FRMGold, anchor=west, fill=white, inner sep=1pt]
    at (\x+0.16,-0.45) {\lab};
\node[font=\scriptsize\itshape, text=FRMGrey, anchor=west] at (0.15,0.42)
  {default assumed to occur mid-period};
\end{tikzpicture}
\end{center}
\figcap{The timing asymmetry that makes the valuation work. Premium payments land
on the anniversary dates; defaults are assumed to land halfway between them.
That is why the payoff and accrual legs are discounted at 0.5, 1.5 and 2.5 while
the payment leg is discounted at 1, 2 and 3, and why an accrual term exists at
all.}

\begin{crexbox}[Valuing a 3-year CDS]
The hazard rate of the reference entity is 3\% per annum, the risk-free rate is
4\% continuously compounded, and the recovery rate is 35\%. Find the breakeven
CDS spread on \$1 of notional.
\tcblower
\textbf{Step 1. Survival and default probabilities.}
\[ \mathrm{PS}_t = e^{-\lambda t} \qquad\qquad
   \mathrm{PD}_t = \mathrm{PS}_{t-1} - \mathrm{PS}_t \]

\centering\footnotesize
\begin{tabular}{c r r}
\toprule
\thead{Year} & \thead{Probability of survival} & \thead{Probability of default} \\
\midrule
1 & 0.97045 & 0.0296 \\
2 & 0.94176 & 0.0287 \\
3 & 0.91393 & 0.0278 \\
\bottomrule
\end{tabular}

\raggedright
\smallskip
\textbf{Step 2. PV of expected payments, per unit of spread $s$.}

\centering\footnotesize
\begin{tabular}{c r r r}
\toprule
\thead{Year} & \thead{Survival} & \thead{Discount factor} & \thead{PV of expected payment} \\
\midrule
1 & 0.97045 & 0.9608 & $0.9324s$ \\
2 & 0.94176 & 0.9231 & $0.8693s$ \\
3 & 0.91393 & 0.8869 & $0.8106s$ \\
\midrule
 & & \textbf{Total} & $\mathbf{2.6123s}$ \\
\bottomrule
\end{tabular}

\raggedright
\smallskip
\textbf{Step 3. PV of expected accrual payments, at the mid-year default dates.}

\centering\footnotesize
\begin{tabular}{c r r r r}
\toprule
\thead{Time} & \thead{PD} & \thead{Accrual} & \thead{Discount factor} &
\thead{PV of expected accrual} \\
\midrule
0.5 & 0.0296 & $0.0148s$ & 0.9802 & $0.0145s$ \\
1.5 & 0.0287 & $0.0143s$ & 0.9418 & $0.0135s$ \\
2.5 & 0.0278 & $0.0139s$ & 0.9048 & $0.0126s$ \\
\midrule
 & & & \textbf{Total} & $\mathbf{0.0406s}$ \\
\bottomrule
\end{tabular}

\raggedright
\smallskip
\textbf{Step 4. PV of expected payoff, with a 35\% recovery rate.}

\centering\footnotesize
\begin{tabular}{c r r r r r}
\toprule
\thead{Time} & \thead{PD} & \thead{RR} & \thead{Expected payoff} &
\thead{Discount factor} & \thead{PV of expected payoff} \\
\midrule
0.5 & 0.0296 & 35\% & 0.0192 & 0.9802 & 0.0188 \\
1.5 & 0.0287 & 35\% & 0.0187 & 0.9418 & 0.0176 \\
2.5 & 0.0278 & 35\% & 0.0181 & 0.9048 & 0.0164 \\
\midrule
 & & & & \textbf{Total} & \textbf{0.0528} \\
\bottomrule
\end{tabular}

\raggedright
\smallskip
\textbf{Step 5. Set the two sides equal.}
\[ \mathrm{PV}\bigl[E(\text{payments})\bigr] = \mathrm{PV}\bigl[E(\text{payoff})\bigr] \]
\[ 2.6529\,s = 0.0528 \]
\[ s = 0.0528 / 2.6529 = 0.0199 \text{, or } 1.99\%\text{, or } 199\text{ bps} \]
\end{crexbox}

\begin{crfmlbox}[The breakeven spread in one line]
\[ s \;=\; \frac{C}{A+B} \]
\vk{$A$ = PV of expected payments per unit of spread; $B$ = PV of expected accrual
payments per unit of spread; $C$ = PV of the expected payoff. In the example above
$A = 2.6123$, $B = 0.0406$ and $C = 0.0528$. Forgetting $B$ raises the computed
spread, because the buyer is then credited with less value paid in.}
\end{crfmlbox}

\subsection*{Marking a CDS to market}

\begin{crfmlbox}[Mark-to-market value]
\[ \text{MtM value to the protection seller} \;=\;
   \mathrm{PV}\bigl[E(\text{payments})\bigr] -
   \mathrm{PV}\bigl[E(\text{payoff})\bigr] \]
\vk{Stated from the \emph{seller's} side: the seller receives the premium leg and
pays the protection leg, so the difference is positive to the seller when spreads
have narrowed since inception. The buyer's mark is the negative of this.}
\end{crfmlbox}

\subsection*{Binary CDS}

A binary CDS is a variation of a regular CDS with one key difference: upon
default, the \term{full notional amount} is paid, regardless of the actual
recovery rate. The price of protection for a binary CDS is therefore typically
\emph{higher} than for a regular CDS.

% =====================================================================
\lo{CR-13 c}{Identify the default probabilities used to value a CDS.}

\begin{crkeybox}[Which probabilities]
Default probabilities used for valuing a CDS should be \term{risk-neutral}
default probabilities. They can be estimated from bond prices, or derived from
CDS quotes.
\end{crkeybox}

This follows the rule set out at CR-9 k: risk-neutral probabilities are for
valuation, real-world probabilities for scenario analysis. A CDS valuation is a
valuation.

% =====================================================================
\lo{CR-13 d}{Evaluate the use of credit indices and fixed coupons in pricing CDS transactions.}

\subsection*{Credit indices}

Credit indices are developed for tracking CDS spreads. Roughly, an index
represents the average CDS spread on the underlying portfolio companies. The key
portfolios are \term{CDX NA IG} (North America) and \term{iTraxx Europe}
(Europe), each covering 125 reference entities.

\begin{crexbox}[Cost of index protection]
The 5-year CDX NA IG index is quoted as ask 66 basis points. A trader wants
\$800,000 of protection on each company. What is the annual cost, and what happens
on a default?
\tcblower
\[ 0.0066 \times \$800{,}000 \times 125 = \$660{,}000 \text{ per year} \]
A company default reduces the annual payment by
\[ \$660{,}000 / 125 = \$5{,}280 \]
\end{crexbox}

\subsection*{Fixed coupons}

The fixed coupon is a standardized fee of 100 basis points (1\%) per year,
typically paid by the protection buyer for investment-grade issuers. It simplifies
contract settlement. Where the CDS spread exceeds the fixed coupon, the buyer pays
an up-front premium. The full three-case table is at \textbf{CR-9 g}.

\begin{crfmlbox}[Price of a CDS, and hence the up-front payment]
\[ \text{price} \;=\; 100 - \bigl[100 \times D \times (s - c)\bigr] \]
\vk{$D$ = the duration of the CDS; $s$ = the CDS spread; $c$ = the fixed coupon,
both as decimals. When $s = c$ the price is 100 and no up-front payment changes
hands. When $s > c$ the price falls below 100 and the \emph{buyer} pays the
difference up front.}
\end{crfmlbox}

% =====================================================================
\lo{CR-13 e}{Define CDS forwards and CDS options.}

\begin{crdefbox}[Forwards and options on a CDS]
A \term{CDS forward} is an \emph{obligation} to buy or sell a particular credit
default swap on a particular reference entity at a particular future time $T$.

\smallskip
A \term{CDS option} is an \emph{option} to buy or sell a particular credit default
swap on a particular reference entity at a particular future time $T$.
\end{crdefbox}

\subsection*{CDS strategies}

\begin{enumerate}
\item If the CDS spread is \emph{less than} the bond yield spread: buy the
corporate bond, and buy CDS protection, to earn more than the risk-free rate.
\item If the CDS spread is \emph{greater than} the bond yield spread: sell the
corporate bond, and sell CDS protection, to borrow at less than the risk-free
rate.
\end{enumerate}

% =====================================================================
\lo{CR-13 f}{Describe the process of valuing a synthetic CDO using the spread payments approach and the Gaussian copula model of time to default approach.}

\begin{crdefbox}[Synthetic CDO]
The issuer typically \emph{sells CDS} on a reference portfolio of underlying
assets. Synthetic CDOs use CDS contracts to mimic the cash flow of a traditional
cash CDO without buying the actual debt. The issuer earns income by selling
protection and pays out if there is a default, just as a traditional CDO does.
\end{crdefbox}

\subsection*{1) Spread payments approach}

\begin{enumerate}[label=\alph*)]
\item Calculate the present value of expected inflows (CDS spread payments) and
outflows (credit losses).
\item Adjust inflows based on the concept of \term{accrued payments}, as in the
CDS valuation at objective b.
\item Set the CDS spread to equate the PV of inflows with the PV of outflows,
achieving a breakeven point: $s = C/(A+B)$.
\end{enumerate}

\subsection*{2) Gaussian copula approach}

\begin{enumerate}[label=\alph*)]
\item Model the default correlations among entities within the CDO using a
Gaussian copula.
\item Compute both unconditional and conditional probabilities of default,
incorporating correlations.
\item Estimate the expected losses for each tranche based on the probability and
severity of defaults.
\end{enumerate}

% =====================================================================
\lo{CR-13 g}{Define the two measures of implied correlation: compound (tranche) correlation and base correlation.}

In the Gaussian copula model the recovery rate is usually assumed to be 40\%.
This leaves the copula correlation $\rho$ as the \emph{only unknown parameter},
so market participants can imply a correlation from market quotes for tranches.

\begin{enumerate}
\item \term{Compound (tranche) correlation.} Unique to each tranche. It is the
correlation that, when plugged into the model, makes the calculated tranche spread
match the market price. It often shows a \term{smile} pattern, with higher
correlation for very risky \emph{and} very safe tranches.
\item \term{Base correlation.} Considers the whole CDO. It adjusts the tranche
correlations from point 1 to ensure the model-derived present value of losses for
each tranche aligns with market prices. This adjustment leads to a \term{skew},
where safer tranches have higher implied correlation.
\end{enumerate}

\begin{center}
\begin{tikzpicture}
\begin{axis}[
  width=12.0cm, height=5.4cm,
  xlabel={Tranche, from most junior to most senior},
  ylabel={Implied correlation},
  xlabel style={font=\sffamily\scriptsize}, ylabel style={font=\sffamily\scriptsize},
  tick label style={font=\sffamily\scriptsize},
  symbolic x coords={equity,junior mezz,mezz,senior,super senior},
  xtick=data, ymin=0, ymax=1, ytick=\empty,
  legend style={font=\scriptsize\sffamily, draw=FRMRule, fill=white,
    at={(0.5,-0.34)}, anchor=north, legend columns=-1, column sep=10pt}]
\addplot[draw=FRMNavy, line width=1.2pt, mark=*, mark size=1.8pt]
  coordinates {(equity,0.62) (junior mezz,0.36) (mezz,0.28) (senior,0.42)
  (super senior,0.68)};
\addlegendentry{compound (tranche) correlation --- a \emph{smile}}
\addplot[draw=FRMGold, line width=1.2pt, mark=square*, mark size=1.8pt]
  coordinates {(equity,0.24) (junior mezz,0.36) (mezz,0.48) (senior,0.62)
  (super senior,0.78)};
\addlegendentry{base correlation --- a \emph{skew}}
\end{axis}
\end{tikzpicture}
\end{center}
\figcap{Illustrative shapes only, no market data. The point is the difference in
\emph{form}: compound correlation is computed tranche by tranche and comes out
high at both ends and low in the middle, which is the smile. Base correlation is
cumulative across the capital structure and rises monotonically, which is the
skew. If asked which measure produces a smile and which a skew, this is the
distinction.}

% =====================================================================
\lo{CR-13 h}{Discuss alternative approaches used to estimate default correlation.}

The standard one-factor Gaussian copula model assumes a single factor drives
defaults, which is not always true. The alternatives are:

\begin{enumerate}
\item \term{Heterogeneous models.} Allow for specification of different
distributions for time to default for different reference entities included in
the collateral pool.
\item \term{Other copulas.} The Student's $t$ copula, Archimedean copula, the
Clayton copula, and the Marshall-Olkin copula.
\item \term{Random recovery rates and factor loadings.} Based on the negative
relationship between default rates and recovery rates, established at CR-9 f.
\item \term{Implied copula model.} Derived from market prices.
\item \term{Dynamic models.} Structural and reduced-form models, which measure the
evolution of loss on a collateral pool over time.
\end{enumerate}
